\subsection{态射的下降}\label{src:3.1.2-subsection}
\begin{proposition}[态射的 fpqc 下降]\label{src:3.1.7}
设 $Y$ 是概形, $f:S'\to S$ 是概形的 fpqc 态射. 若 $g:S'\to Y$ 是满足 $g\circ p_1=g\circ p_2$ 的态射, 则存在唯一的态射 $h:S\to Y$, 将下图补成交换图:
\[
\begin{tikzcd}[column sep=large]
S'\times_S S'\arrow[r,shift left=1.1ex,"p_1"]
\arrow[r,shift right=1.1ex,"p_2"']&
S'\arrow[r,"f"]\arrow[dr,"g"']&S\arrow[d,dashed,"h"]\\
&&Y
\end{tikzcd}
\]
\end{proposition}
换言之, fpqc 态射 $f:S'\to S$ 是概形范畴中的\newterm{有效满态射}{effective epimorphism}, 即对每个概形 $Y$, 序列
\begin{equation}\label{src:3.1.8}
\Mor(S,Y)\longrightarrow\Mor(S',Y)\rightrightarrows\Mor(S'\times_S S',Y)
\end{equation}
正合; 仅为满态射则相当于第一个映射是单射. 我们很快将看到, 这也可以表述为函子 $\Mor(-,Y)$ 是 fpqc 拓扑中的层 (见命题~\bookref{3.3.8}).
\begin{proof}
仿射情形很直接. 写成 $S'=\Spec A'$, $S=\Spec A$ 与 $Y=\Spec R$, 则由命题~\bookref{3.1.1}~知 $A\to A'\rightrightarrows A'\otimes_A A'$ 正合. 应用 $\Hom(R,-)$, 可知 $\Hom(R,A)\to\Hom(R,A')\rightrightarrows\Hom(R,A'\otimes_A A')$ 也正合. 在仿射概形与环的对偶关系下, 这正是~\bookeqref{3.1.8}~的正合性. 为归约到仿射情形, 注意问题在 $S$ 上是局部的, 因而可设 $S$ 仿射. 由于 $S'\to S$ 是 fpqc 的, 存在映满 $S$ 的拟紧开子集 $U'\subseteq S'$. 选取 $U'$ 的有限仿射覆盖 $\{U_i'\}$ 后, 可将 $S'$ 替换为仿射概形 $\coprod_i U_i'$ (细节留给读者).

将 $Y$ 归约到仿射情形稍难一些. 首先注意,
\begin{equation}\label{src:3.1.9}
|S'\times_S S'|\rightrightarrows|S'|\longrightarrow|S|
\end{equation}
是集合的余等化子图. 事实上, 我们已经知道 $|S'|\to|S|$ 满射, 而且若 $s_1',s_2'\in|S'|$ 在 $S$ 中具有相同的像, 则存在 $q\in|S'\times_S S'|$ 满足 $p_1(q)=s_1'$ 与 $p_2(q)=s_2'$. 因此存在集合映射 $|h|:|S|\to|Y|$, 使得 $|g|=|h|\circ|f|$. 设 $U\subseteq Y$ 是仿射开子集, 则 $V:=|h|^{-1}(U)$ 是 $S$ 的子集, 且 $f^{-1}(V)=g^{-1}(U)$ 是 $S'$ 的开子集. 由习题~\bookref{A.4.9}, $S'\to S$ 是浸没映射 (即 $S$ 带有商拓扑), 因而 $V\subseteq S$ 是开的. 又因 $f^{-1}(V)\to V$ 也是 fpqc 的, 我们可以假设 $Y$ 仿射. 另见~\cite[定理 $2$, 第 317 页]{FGAI}, \cite[定理 $6$(a)]{BLR90}, \cite[定理 $2.55$]{Vis05}~及~\cite[Tag 023Q]{SP}.
\end{proof}
下述推广也成立.
\begin{corollary}\label{src:3.1.10}
设 $f:S'\to S$ 是概形的 fpqc 态射.
\begin{enumerate}
\item 若 $S\to T$ 是概形态射, $Y$ 是 $T$-概形, 则
\[
\Mor_T(S,Y)\longrightarrow\Mor_T(S',Y)
\rightrightarrows\Mor_T(S'\times_S S',Y)
\]
正合.
\item 若 $X$ 与 $Y$ 是 $S$ 上的概形, 则
\[
\Mor_S(X,Y)\longrightarrow\Mor_{S'}(X_{S'},Y_{S'})
\rightrightarrows\Mor_{S''}(X_{S''},Y_{S''})
\]
正合, 其中 $S''=S'\times_S S'$.
\end{enumerate}
\end{corollary}
\begin{proof}
对于 (1), 由命题~\bookref{3.1.7}~知 $\Mor_T(S,Y)\to\Mor_T(S',Y)$ 是单射, 而且若 $g:S'\to Y$ 是满足 $p_1\circ g=p_2\circ g$ 的 $T$-态射, 则存在概形映射 $h:S\to Y$ 使 $g=h\circ f$. 记结构态射为 $p_S:S\to T$, $p_{S'}:S'\to T$ 及 $p_Y:Y\to T$. 注意, $p_Y\circ h$ 与 $p_S$ 是 $\Mor(S,T)$ 中的元素, 二者都映到 $\Mor(S',T)$ 中的 $p_Y\circ g=p_{S'}$. 由命题~\bookref{3.1.7}, 包含映射 $\Mor(S,T)\to\Mor(S',T)$ 是单射, 故 $h:S\to Y$ 是 $T$-态射. 将 (1) 应用于 $T=S$ 上的 fpqc 态射 $f:X_{S'}\to X$, 即得 (2).
\end{proof}
